% Part: sets-functions-relations
% Chapter: arithmetization
% Section: reflections

\documentclass[../../../include/open-logic-section]{subfiles}

\begin{document}

\olfileid{sfr}{arith}{ref}
\olsection{Een paar filosofische gedachten}
Tot zover de technische details. Maar wat heeft dat alles opgeleverd?
Eén ding staat vrijwel vast: dit leverde {mooie} zuivere wiskunde op.
Ook kwamen er diepgaande inzichten uit. De volledigheidseigenschap
laat zien wat het doorslaggevende verschil tussen reële en rationale
getallen is. Dat was een diep inzicht. Door precies te beschrijven hoe je
reële getallen als Dedekindsneden opbouwt, krijgt de analyse bovendien een stevige
basis. We weten nu dat een \emph{volledig geordend lichaam} echt kan
bestaan: een getalsysteem met de vereiste rekenregels, volgorde en
volledigheid. De sneden vormen zo'n lichaam.
Toch zijn er een paar redenen om voorzichtig te zijn met deze resultaten.
Ten eerste is het niet duidelijk dat reële getallen als sneden zien
\emph{strenger} is dan ze zien als hun bekende, misschien oneindige
rijen cijfers achter de komma. Die laatste ``opbouw'' lijkt wat op de
opbouw met Cauchyrijen. Toch kenden wiskundigen het idee in wezen al
vanaf het begin van de zeventiende eeuw (zie \olref[cauchy]{sec}). De
echte winst aan strengheid kwam door het inzicht dat de reële getallen de
volledigheidseigenschap hebben. Dat we ze ook als bepaalde
verzamelingen kunnen opbouwen, is op zichzelf misschien niet zo
interessant.
Nog minder duidelijk is of het veel makkelijker opbouwen van de gehele en
rationale getallen binnen de verzamelingenleer het wiskundige werk strenger maakt. Let op wat we in de
praktijk doen. Eerst hebben we de gehele getallen als klassen van
gelijkwaardige geordende paren natuurlijke getallen
\emph{opgebouwd}. Daarna bouwden we de rationale getallen op als zulke
klassen van paren gehele getallen, en de reële getallen als
verzamelingen rationale getallen. Meteen daarna gaan we die hele opbouw
weer \emph{vergeten}. Niemand wil hem tijdens een wiskundig bewijs
\emph{gebruiken}, behalve natuurlijk om te bewijzen dat de opbouw zelf
goed werkt. Rechtstreeks over een reëel getal praten is veel
makkelijker dan praten over een verzameling van verzamelingen van
verzamelingen van verzamelingen van verzamelingen van verzamelingen van
verzamelingen natuurlijke getallen.
%(For reals were constructed as sets of rationals; and these were constructed as equivalence classes of ordered pairs of integers, i.e., sets of sets of sets of integers; etc.).
%2/3 = [2,3]
%	i.e. {<2,3>, ... }
%	{ {{2},{2,3}}, ... }
%
%reals are 
%	sets of rationals 
%
%rationals are 
%	sets of sets of sets of integers
%
%integers are
%	sets of sets of sets of naturals
Het minst geloofwaardig is dat deze definities ons vertellen wat
gehele, rationale en reële getallen \emph{zijn}, \emph{als we vragen
wat ze werkelijk zijn}. Het is dus maar de vraag of reële getallen
bijvoorbeeld \emph{niets anders zijn dan} bepaalde verzamelingen van
verzamelingen van verzamelingen\ldots. Het grootste probleem voor dat idee
is dat we de opbouw op veel andere manieren hadden kunnen doen. Voor de reële getallen bestaan echt
interessante andere manieren (zie \olref[cauchy]{sec}). Maar het kan
ook heel eenvoudig anders. In \olref[sfr][rel][ref]{sec} zagen we dat
we een geordend paar iets anders hadden kunnen definiëren. Met die
andere definitie had alles even goed gewerkt, maar waren er andere
objecten uitgekomen. Er zijn dus meerdere manieren om de gehele,
rationale en reële getallen als verzamelingen op te bouwen. Dan zou het
zomaar een keuze zijn, en nogal gênant, om te beweren dat de gehele
getallen bijvoorbeeld \emph{deze} verzamelingen zijn in plaats van
\emph{die}. (Zoals in \olref[sfr][rel][ref]{sec} is dit een voorbeeld
van een argument dat bekend is geworden door
\citealt{Benacerraf1965}.)
Er is nog iets \emph{vreemds} aan onze opbouw. We begonnen met de
natuurlijke getallen. Daarna bouwden we de gehele getallen, met daarin
``de $0$ van de gehele getallen'': $ \equivrep{0,0}{\Intequiv}$. Maar
$0 \neq \equivrep{0,0}{\Intequiv}$. Volgens deze opbouw is
\emph{geen enkel} natuurlijk getal een geheel getal. Dat gaat sterk
tegen ons gewone idee van getallen in. In
\olref[sfr][set][imp]{sec} zeiden we immers bijna zonder uitleg dat
$\Nat \subseteq \Rat$. Als de opbouw precies vertelt \emph{wat} de
getallen zijn, was meteen duidelijk dat dit niet klopte.
Wat houden we dan over? In een na\"ieve verzamelingenleer kunnen we de
natuurlijke getallen als gegeven nemen en de gehele, rationale en reële
getallen als bepaalde verzamelingen \emph{behandelen}. Zo kunnen we de
theorieën over deze getallen in de verzamelingenleer \emph{onderbrengen}.
Maar wat dat filosofisch betekent, is niet zo eenvoudig.
Dit is natuurlijk niet het laatste woord!{} Het punt is alleen dit. Wie
wil laten zien dat het opbouwen van de reële getallen binnen de
verzamelingenleer filosofisch echt
belangrijk \emph{is}, heeft daar nog een \emph{filosofisch} argument
voor nodig.
%Additionally: for the entire duration of this chapter, we have treated
%the natural numbers as simply \emph{given} to us. But what can we do
%with them? That is the subject matter for the next chapter.
\end{document}
